\chapter{Thermodynamics of Bodies}

Up to this point: \emph{purely mechanical} primitive notions.\\

Introduce new concepts:
%\begin{enumerate}[label=(\Alph*)]

\begin{enumerate}
\item 
Internal energy (specific internal energy)

\item
Heating (heat flux and heat supply)

\end{enumerate}

Those will be restricted together with other constituents of the early discussion of motions (dynamic processes) to be obey the 

\begin{enumerate}[label=(\Alph*)]

\item
Axiom of energy conservation (the first law of thermodynamics)\\
further introduce notions of 

\begin{enumerate}[label=\arabic*.]
\setcounter{enumii}{2}
\item
Entropy (or Carlory) $\cdots$ specific entropy

\item
Absolute temperature

\end{enumerate}
and all of the above will obey a second postulate called the 

\item
Entropy inequality (Clausius-Duhem or Truesdell-Toupin) which is a particular version of the second law of thermodynamics.

\end{enumerate}

\section{Energetic Axioms}

Let $\tenq{y} = \hat{\tenq{y}}(\tenq{x},t)$ be and admissible motions. \\[3pt]


\fbox{Principle 1.}\\[3pt]

There exists a scalar field $\varepsilon(\bullet,t)\colon R_t \to R\phantom{-}\forall t \in \tau$ with $\varepsilon\in C^1(R_t \times \tau)$ called the specific internal energy (energy per unit mass) such that 

\[\int_{P_t}\varepsilon(\tenq{y},t)\rho(\tenq{y},t)d\volume_{\tenq{y}} = E(P_t)\]

\[E(P_t)\cdots\text{total internal energy of }P_t \subset R_t\]
dim $\varepsilon = $ Energy/mass\\[3pt]


\fbox{Principle 2.}\\[3pt]

There exists for each $\tenq{n}\in u$ a scalar field $\hslash(\tenq{n},\bullet,t)\in C^1(R_t)$ called \emph{the heat flux} such that

\begin{align*}
\int_{\partial P_t}\hslash(\tenq{n}(\tenq{y},t),\tenq{y},t)dA_{\tenq{y}} 
&= \text{total contact heating}\\
&\equiv Q^c \text{ (due to heat conduction)}
\end{align*}

dim $\hslash = \frac{\text{Energy}}{\text{(unit area)(unit time)}}$\\
dim $Q^c = \frac{\text{Energy}}{\text{unit time}}$\\[3pt]

\fbox{Principle 3.}\\[3pt]

There exists a scalar field $r(\bullet,t)\in C(R_t)$ called the specific heat supply (per unit mass) such that 


\begin{align*}
\int_{P_t}r(\tenq{y},t)\rho(\tenq{y},t)d\volume_{\tenq{y}}
&= \text{Body heating}\\
&\equiv Q^r \text{ (due to radiation)}
\end{align*}




Total heating of $P_t$

\begin{align*}
Q(P_t) = Q^c + Q^r 
&= \int_{\partial P_t}\hslash(\tenq{n},\tenq{y},t)dA_{\tenq{y}}\\
&+ \int_{P_t}r(\tenq{y},t)\rho(\tenq{y},t)d\volume_{\tenq{y}}
\end{align*}


For a motion recall also that

\[K(P_t) = \int_{P_t} \frac{1}{2}\rho(\tenq{y},t)\abs{\bar{\tenq{v}}(\tenq{y},t)}^2 d\volume_{\tenq{y}}\]

While the mechanical power $P(P_t)$

\begin{align*}
P(P_t) 
&= \int_{\partial P_t}\tenq{t}(\tenq{n}(\tenq{y},t),\tenq{y},t)\cdot\bar{\tenq{v}}(\tenq{y},t)dA_{\tenq{y}}\\
&+ \int_{P_t}\tenq{b}(\tenq{y},t)\cdot\bar{\tenq{v}}(\tenq{y},t)d\volume_{\tenq{y}}
\end{align*}

\begin{axiom}{Conservation (Balance) of Energy}

$\forall P_t \in R_t$\\

\[\frac{d}{dt}\left[E(P_t) + K(P_t)\right] = P(P_t) + Q(P_t)\]

Namely:\\

The rate of change of internal + kinetic energy  of each subbody equals the mechanical power + the heating. 

\begin{equation}
\begin{split}
\Rightarrow 
&\frac{d}{dt}\int_{P_t}\rho\left[\varepsilon + \frac{1}{2}\bar{\tenq{v}}\cdot\bar{\tenq{v}}\right]d\volume\\
&= \int_{P_t}\left[\tenq{b}\cdot\bar{\tenq{v}} + r\rho\right]d\volume + \int_{\partial P_t}\left[\hslash(\tenq{n}) + \tenq{t}(\tenq{n})\cdot\bar{\tenq{v}}\right]dA
\end{split}
\tag{$\circledS$} \label{Energetic Axioms_circles} 
\end{equation}
\end{axiom}

\begin{theorem}{Local Energy Balance, Heat flux vector}

Let $\hat{\tenq{y}}\colon R\to R_t$ with mass density $\rho(\bullet,t)$ and spatial velocity $\bar{\tenq{v}}(\bullet,t)$ on $R_t$. Then the specific internal energy $\varepsilon(\bullet,t)$, the heat flux $\hslash(\tenq{n},\bullet,t)$ and heat supply $r(\bullet,t)$ on $R_t$ abiding by Principle 1, Principle 2, and Principle 3, are such that $\eqref{Energetic Axioms_circles}$ holds \emph{if and only if} 

\begin{enumerate}[label=(\roman*)]

\item
There exists a vector field 

\[\tenq{q}(\bullet,t)\colon R_t\to E_3, \tenq{q}(\bullet,t)\in C^1(R_t)\]

such that

\[\hslash(\tenq{n},\tenq{y},t) = -\tenq{q}(\tenq{y},t)\cdot\tenq{n}\phantom{-}\forall\tenq{y}\in R_t\phantom{-}\forall\tenq{n}\in u\]

\item
\begin{align*}
&\tenq{\tau}(\tenq{y},t)\cdot\tenq{D}(\tenq{y},t)-\nabla\cdot\tenq{q}(\tenq{y},t) + \rho(\tenq{y},t)r(\tenq{y},t)\\
&= \rho(\tenq{y},t)\dot{\varepsilon}(\tenq{y},t)\phantom{-}\forall (\tenq{y},t)\in R_t\times\tau
\end{align*}

\end{enumerate}

\end{theorem}

\begin{proof}\mbox{}\\

First recall the power identity: 

\begin{equation}
\frac{d}{dt}K(P_t) = P(P_t) - \int_{P_t}\tenq{\tau}\cdot\tenq{D}d\volume
\tag{1} \label{Energetic Axioms-(1)}
\end{equation}

Also 

\begin{equation}
\frac{d}{dt}E(P_t) = \frac{d}{dt}\int_{P_t}\rho\varepsilon d\volume \vc{=}{RTT}{} \int_{P_t}\rho(\tenq{y},t)\dot{\varepsilon}(\tenq{y},t)d\volume
\tag{2} \label{Energetic Axioms-(2)}
\end{equation}

where $\dot{\varepsilon} = \frac{\partial \varepsilon}{\partial t} + \left(\nabla\varepsilon\right)\cdot\bar{\tenq{v}}$\\

Substitute \eqref{Energetic Axioms-(1)} and \eqref{Energetic Axioms-(2)} in $\eqref{Energetic Axioms_circles}$

\begin{equation}
\int_{P_t}\rho\dot{\varepsilon}d\volume = \int_{\partial P_t}\hslash(\tenq{n})dA + \int_{P_t}\left[\rho r + \tenq{\tau}\cdot\tenq{D}\right]d\volume
\tag{3} \label{Energetic Axioms-(3)}
\end{equation}

Repeat proof used in obtaining $\tenq{t} = \tenq{\tau}\tenq{n}$, i.e. choose $P_t$ to be a tetrahedron $P_t(h)$, with height $h$ and let $h\to 0$. 

\begin{equation}
\Rightarrow \hslash(\tenq{n},\tenq{y},t) = -\tenq{q}(\tenq{y},t)\cdot\tenq{n}
\tag{4} \label{Energetic Axioms-(4)}
\end{equation}

\[\eqref{Energetic Axioms-(3)} \& \eqref{Energetic Axioms-(4)} \Rightarrow\int_{P_t}\rho\dot{\varepsilon}d\volume = -\int_{\partial P_t}\tenq{q}\cdot\tenq{n}dA + \int_{P_t}\left[\rho r + \tenq{\tau}\cdot\tenq{D}\right]d\volume\]

\[\vc{\Rightarrow}{divegence }{theorem} \int_{P_t} \left[\rho\dot{\varepsilon} + \nabla\cdot\tenq{q}-\rho r -\tenq{\tau}\cdot\tenq{D}\right]d\volume = 0\]

Smoothness of $\rho$, $\varepsilon$, $\hslash$ (or $\tenq{q}$), $r$, $\cdots$ and localization argument gives

\[\tenq{\tau}\cdot\tenq{D}-\nabla\cdot\tenq{q}+\rho r = \rho \dot{\varepsilon}\]

\end{proof}

Consequences of the \\

\fbox{First law}
\[\int_{P_t}\rho\dot{\varepsilon}d\volume = \int_{\partial P_t}\underbrace{\hslash(\tenq{n})}_{\text{heat flux}}dA + \int_{P_t}\left[\underbrace{\rho r}_{\text{specific heat supply}} + \tenq{\tau}\cdot\tenq{D}\right]d\volume\]

\[\Rightarrow \hslash(\tenq{n},\tenq{y},t) = -\underbrace{\tenq{q}(\tenq{y},t)}_{\text{heat flux vector}}\cdot\tenq{n}\]

also


\begin{equation}
\boxed{
\tenq{\tau}\cdot\tenq{D}-\nabla\cdot\tenq{q}+\rho r = \rho \dot{\varepsilon}
}\cdots\text{local version of $1^{st}$ law}
\tag{$\ast$} \label{Energetic Axioms-(*)}
\end{equation}

\section{Motivation and comparison with classical thermodynamics}

A thermodynamics consider processes which are homogeneous (independent of position $\tenq{y}$). 


\begin{alignat*}{3}
& \text{Internal energy }     \quad && \colon \quad && \varepsilon(t)\\
& \text{Pressure } \quad && \colon 	\quad &&  p(t)\\
& \text{Volume }  \quad && \colon	\quad && v(t)\\
& \text{Heating } \quad && \colon	\quad && \tilde{q}(t)
\end{alignat*}
%& \phantom{-}  \quad && \phantom{-}	\quad && \phantom{-}	\quad && \text{(rotation)} \\	
(no distinction between conduction and radiation)\\

The first law (compare to $\eqref{Energetic Axioms-(*)}$)

\[\dot{\varepsilon} = p\dot{v} + \tilde{q}\]
\begin{align*}
\dot{\varepsilon}
&\cdots \rho\dot{\varepsilon}(\tenq{y},t)\\
p\dot{v}
&\cdots \tenq{\tau}(\tenq{y},t)\cdot\tenq{D}(\tenq{y},t)\\
\tilde{q}
&\cdots -\nabla\cdot\tenq{q}(\tenq{y},t) + \rho(\tenq{y},t)r(\tenq{y},t)
\end{align*}

For the second law one introduces an entropy $\eta(t)$ and an absolute temperature $\theta(t) > 0$. The second law assumes an upper bound for the quantity $\tilde{q}(t)/\theta(t)$, i.e. 

\[\dot{\eta}(t)\geq \frac{\tilde{q}(t)}{\theta(t)}\]

(Entropy rate exceeds or equals heating over temperature)\\

First and second law 
\[\Rightarrow \theta \dot{\eta}-\left(\dot{\varepsilon}-p\dot{v}\right)\geq 0\]

(Entropy rate $\times$ temperature $\geq$ internal energy rate $-$ mechanical power \footnote{rate of energy \emph{not} produced by mechanical work which is reversible process})


\section{Axioms of entropy and temperature}


\fbox{Principle 4.}\\
Given a motion $\hat{\tenq{y}}(\bullet,t)\colon R\to R_t$, $t \in [t_0,t_1]=\tau$, define a scalar field $\eta(\bullet,t)\colon R_t\to \mathbb{R}$ called the specific entropy $\eta\in C^1(R_t)$ such that 

\[\int_{P_t}\rho(\tenq{y},t)\eta(\tenq{y},t)d\volume_{\tenq{y}} = S(P_t)\cdots\text{total entropy}, P_t\subset R_t\]

A strictly positive scalar field $\theta(\bullet,t)\colon R_t \to R$ ($\theta(\tenq{y},t)>0\phantom{-}\forall \tenq{y}\in R_t \phantom{-}\forall t\in\tau$), $\theta\in C^1(R_t)$ called the absolute temperature field. The following axiom is commonly accepted as a statement of the second law in continuous  mechanics. 

\[\frac{dS(P_t)}{dt}\geq-\int_{\partial P_t}\frac{1}{\theta(\tenq{y},t)}\tenq{q}(\tenq{y},t)\cdot\tenq{n}(\tenq{y},t)dA_{\tenq{y}} + \int_{P_t}\frac{\rho(\tenq{y},t)r(\tenq{y},t)}{\theta(\tenq{y},t)}d\volume_{\tenq{y}}\]

where

\[S(P_t) = \int_{P_t}\rho(\tenq{y},t)\eta(\tenq{y},t)d\volume_{\tenq{y}}\]

\begin{remark}\mbox{}\\

Call the entropy production rate of $P_t$

\[\Gamma(P_t)\equiv\frac{d}{dt}\int_{P_t}\rho\eta d\volume + \int_{\partial P_t}\frac{\tenq{q}\cdot\tenq{n}}{\theta}dA - \int_{P_t}\frac{\rho r}{\theta}d\volume\]

Thus second law (Clausius-Duhem)

\[\Leftrightarrow \Gamma(P_t)\geq 0\phantom{-}P_t\subset R_t\phantom{-}t \in \tau\]

recall that $\frac{d}{dt}\int_{P_t}\rho\eta d\volume = \int_{P_t}\rho\dot{\eta}d\volume$ by RTT

where $\dot{\eta} = \frac{\partial\eta}{\partial t} + \nabla\eta\cdot\bar{\tenq{v}}$

by using the divergence theorem and the above 

\[\Rightarrow \rho\dot{\eta} + \nabla\cdot\left(\frac{\tenq{q}}{\theta}\right) - \frac{\rho r}{\theta}\geq 0\]

\begin{equation}
\Rightarrow
\boxed{
\rho\dot{\eta} + \frac{1}{\theta}\nabla\cdot\tenq{q}-\frac{1}{\theta^2}\tenq{q}\cdot\nabla\theta - \frac{\rho r}{\theta}\geq 0
}
\tag{$\ast\ast$}\label{Energetic Axioms-(**)}
\end{equation}

recall $\eqref{Energetic Axioms-(*)}$ (local version of the first law, $\nabla\cdot\tenq{q} = -\rho\dot{\varepsilon} + \tenq{\tau}\cdot\tenq{D} + \rho r$)

\[\eqref{Energetic Axioms-(*)},\eqref{Energetic Axioms-(**)}\Rightarrow 
\boxed{
\rho\theta\dot{\eta} - \rho\dot{\varepsilon} + \tenq{\tau}\cdot\tenq{D} - \frac{1}{\theta}\tenq{q}\cdot\nabla\theta\geq 0
}
\]
\[\cdots\text{reduced dissipation inequality}\]


\end{remark}

\begin{theorem}{(Dissipation Inequalities)}{(Local)}

The Clausius-Duhem inequality is equivalent to 

\begin{enumerate}[label=(\roman*)]

\item
local entropy production inequality

\begin{equation}
\rho\dot{\eta} + \nabla\cdot\left(\frac{1}{\theta}\tenq{q}\right) - \frac{\rho r}{\theta}\geq 0
\tag{i}\label{Energetic Axioms-(i)}
\end{equation}

\item
reduced dissipation inequality (the first law $+$ the second law)

\begin{equation}
\rho\theta\dot{\eta} - \left(\rho\dot{\varepsilon} - \tenq{\tau}\cdot\tenq{D}\right) - \frac{1}{\theta}\tenq{q}\cdot\nabla\theta\geq 0
\tag{ii}\label{Energetic Axioms-(ii)}
\end{equation}

\end{enumerate}

\begin{note}

equivalent in classical thermodynamics

\[\theta\dot{\eta}-\left(\dot{\varepsilon}-p\dot{v}\right)\geq 0\]
(i.e., $\nabla\theta = 0$, homogeneous assumption)
\end{note}

\end{theorem}

\begin{remark}\mbox{}\\

\begin{itemize}

\item
The continuum and classical versions of dissipation inequalities are equivalent as expected.

\item
The term $\frac{1}{\theta}\tenq{q}\cdot\nabla\theta$ is missed in classical thermodynamics because of the homogeneity assumption. 


\end{itemize}

Thermodynamics deal with \emph{equilibrium system} only.


\end{remark}


\emph{Define} the \emph{internal dissipation} rate $\delta$ as 

\[\delta = \theta\dot{\eta} - \left[\dot{\varepsilon}-\frac{1}{\rho}\tenq{\tau}\cdot\tenq{D}\right]\]

and the \emph{conductive entropy production} density $\xi$ as 

\[\xi = -\frac{1}{\theta}\tenq{q}\cdot\nabla\theta\]

Entropy production

\[\Gamma(P_t) = \int_{P_t}\frac{\rho(\tenq{y},t)\delta(\tenq{y},t) + \xi(\tenq{y},t)}{\theta(\tenq{y},t)}d\volume_{\tenq{y}}\geq 0\]

is the sum of the internally dissipation and heat conduction pants. The inequalities $\delta\geq 0$ and $\xi\geq 0$ are known as the Clausius-Plank and the Fourier Inequality (in heat transfer). We do not view each of them as valid separately in a general process where both mechanical and conductive effects occurs. These inequalities are not implied by the second law. \\[5pt]



Specific \emph{Helmholtz free energy} $\psi(\bullet,t)\colon R_t\to \mathbb{R}$\\

define

\[\psi(\tenq{y},t) = \varepsilon(\tenq{y},t)-\theta(\tenq{y},t)\eta(\tenq{y},t)\phantom{-}\forall\tenq{y}\in R_t\phantom{-}\forall t\in \tau\]

Take Lagrangian derivative

\begin{equation}
\dot{\psi} = \dot{\varepsilon} - \dot{\theta}\eta - \theta\dot{\eta}
\tag{$\star$}\label{Energetic Axioms-(star)}
\end{equation}

\eqref{Energetic Axioms-(i)}, \eqref{Energetic Axioms-(ii)}, and $\eqref{Energetic Axioms-(star)} \Rightarrow$

\[-\rho\left(\dot{\psi}+\eta\dot{\theta}\right) + \tenq{\tau}\cdot\tenq{D}-\frac{1}{\theta}\tenq{q}\cdot\nabla\theta \geq 0\]

\begin{remark}\mbox{}\\

\begin{figure}[H]
\centering
\resizebox{0.2\textwidth}{!}{\input{figure33.eps_tex}}
\caption{Scheme for $P_t$}
\label{Scheme for $P_t$}
\end{figure}

\begin{enumerate}[label=(\arabic*)]

\item
We say that $P_t$ is thermally isolated if $\tenq{q}\cdot\tenq{n} = 0$ on $\partial P_t$ and $r=0$ on $P_t$ (closed) system.

\[\frac{d}{dt}S(P_t)\geq 0\phantom{---}(\text{in classical thermodynamics, }\dot{\eta}\geq\frac{\tilde{q}}{\theta}=0)\]

i.e., entropy cannot decrease under thermal isolation.

\item
Special cases

\begin{align*}
\dot{\varepsilon} = \dot{\eta} = 0
&&
\text{or}
&&
\dot{\psi} = \dot{\theta} = 0
&&
\rightarrow \text{steady state}
\end{align*}

\[\Rightarrow \tenq{\tau}\cdot\tenq{D} - \frac{1}{\theta}\tenq{q}\cdot\nabla\theta\geq 0\]

\begin{enumerate}[label=case (\roman*)]

\item
Non-conductive material ($\tenq{q}=\tenq{0}$ or $\theta$ is uniform, $\nabla\theta = 0$)

\[\Rightarrow\tenq{\tau}\cdot\tenq{D}\geq 0\text{ (non-negetive stress-power density)}\]

\item
For a rest motion $\tenq{D} = \tenq{0}$

\[\Rightarrow -\frac{1}{\theta}\tenq{q}\cdot\nabla\theta\geq 0\Rightarrow\tenq{q}\cdot\nabla\theta\leq 0\cdots \text{Fourier's inequality}\]


\end{enumerate}

\item
If $\tenq{q} = \tenq{0}$ (adiabatic)

\begin{align*}
\delta
&= \theta\dot{\eta} - \left(\dot{\varepsilon} - \frac{1}{\rho}\tenq{\tau}\cdot\tenq{D}\right)\geq 0\\
&= -\left(\dot{\psi} + \eta\dot{\theta}\right)+\frac{1}{\rho}\tenq{\tau}\cdot\tenq{D}\geq 0
\end{align*}


\begin{alignat*}{3}
& \text{If the process is \emph{isentropic} }(\dot{\eta} = 0)     \quad && \Rightarrow \quad && \frac{1}{\rho}\tenq{\tau}\cdot\tenq{D}\geq\dot{\varepsilon}\\
& \text{If the process is \emph{isothermal} }(\dot{\theta} = 0)     \quad && \Rightarrow \quad && \frac{1}{\rho}\tenq{\tau}\cdot\tenq{D}\geq\dot{\psi}
\end{alignat*}



\end{enumerate}


\end{remark}

\section{Thermodynamic process}

\begin{figure}[H]
\centering
\resizebox{0.5\textwidth}{!}{\input{figure34.eps_tex}}
\caption{Scheme for $\tenq{\chi}_t$}
\label{Scheme for chi}
\end{figure}

\begin{definition}\mbox{}\\

Let $\tenq{\chi}_t\colon B\to E_3$, $t\in\tau(\tenq{\chi}_t\in\Omega(B))$ be a motion of body $B$ with $R_t = \tenq{\chi}_t(B)$. The ordered array of 8 elements $\left[\tenq{\chi}_t, \tenq{\tau}, \tenq{b}, \varepsilon, \tenq{q}, r, \eta, \theta\right]$ is a thermodynamic process undergone by $B$ during $\tau$ if $\left[\tenq{\chi}_t, \tenq{\tau}, \tenq{b}\right]$ is a dynamical process and $\varepsilon, \tenq{q}, r, \eta, \theta$ satisfy balance of energy and the entropy inequality in which $\rho$ is governed by mass balance.

\end{definition}

\section{Fundamental field equations for a thermodynamic process}

Let $\bar{\tenq{v}}(\bullet,t)\colon R_t\to E_3$ is spatial (Eulerian) velocity field due to the motion and assume that the density $\rho\in C^1(\mathfrak{m})$, $\mathfrak{m} = \left\{(\tenq{y},t)\mid \tenq{y}\in R_t, t\in\tau\right\}$ then $\left[\tenq{\chi}_t, \tenq{\tau}, \tenq{b}, \varepsilon, \tenq{q}, r, \eta, \theta\right]$ is a thermodynamic process \emph{if and only if}:

\begin{enumerate}[label=(\arabic*)]
\setcounter{enumi}{-1}

\item\label{Fundamental field equations for a thermodynamic process-item-0}
$\dot{\rho} + \rho\nabla\cdot\bar{\tenq{v}}=0\cdots$(mass balance)

\item\label{Fundamental field equations for a thermodynamic process-item-1}
$\nabla\cdot\tenq{\tau} + \tenq{b} = \rho\bar{\tenq{a}}\cdots$(linear momentum balance)

\item
$\tenq{\tau} = \tenq{\tau}^T\cdots$(angular momentum balance)

\item
$\rho\dot{\varepsilon} = -\nabla\cdot\tenq{q} + \tenq{\tau}\cdot\tenq{D} + \rho r\cdots$(the first law: energy balance)

\item \label{Fundamental field equations for a thermodynamic process-item-4}
$\rho\dot{\eta} + \nabla\cdot\left(\frac{1}{\theta}\tenq{q}\right) - \frac{\rho r}{\theta}\geq 0\cdots$(the second law: energy inequality)\\

in the above \ref{Fundamental field equations for a thermodynamic process-item-4} is interchangeable with 
\begin{enumerate}[label=(\roman*)]

\item
$\rho\theta\dot{\eta} - \rho\dot{\varepsilon} + \tenq{\tau}\cdot\tenq{D} - \frac{1}{\theta}\tenq{q}\cdot\nabla\theta\geq 0$

\item
$-\rho\left(\dot{\psi} + \eta\dot{\theta}\right) + \tenq{\tau}\cdot\tenq{D}-\frac{1}{\theta}\tenq{q}\cdot\nabla\theta\geq 0$

\end{enumerate}
\end{enumerate}

where $\bar{\tenq{a}} = \dot{\bar{v}}$, $\tenq{D} = \textit{sym }\nabla\bar{\tenq{v}}$ on $\mathfrak{m}$\\

\ref{Fundamental field equations for a thermodynamic process-item-0}-\ref{Fundamental field equations for a thermodynamic process-item-4} are the fundamental field equations for thermodynamic process. 


\begin{remark}\mbox{}\\
The above are valid for all bodies and admissible motions. Given a \emph{motion} and a \emph{temperature field}, that exists infinitely many choice of $\tenq{\tau}$, $\varepsilon$, $\tenq{q}$, $\eta$, $\tenq{b}$, and $r$ such that \ref{Fundamental field equations for a thermodynamic process-item-1}-\ref{Fundamental field equations for a thermodynamic process-item-4} hold. Moreover for all different types of bodies subjected to the same motion and temperature fields, in general, will undergo different thermodynamic processes. Therefore, the constitutive laws are needed.
\end{remark}

\begin{note}\mbox{}\\
We assumed that $\rho$ obeys mass balance. 
\end{note}